\documentclass[12pt]{article}
\usepackage{kyradefaults}
\addbibresource{references.bib}

\title{\textbf{All Aboard? Causal Evidence on the Labor-Market Effects of New Rail Stations}}
\author{Kyra Sadovi%
  \thanks{E-Mail: \texttt{ksadovi@uchicago.edu}}}
\affil{Harris School of Public Policy, University of Chicago}
\date{}

\begin{document}
\maketitle
\thispagestyle{empty}
 
Public transit is among the most consequential infrastructural interventions a city can make, yet credible causal evidence on its effects on workers' economic outcomes remains limited. This paper asks: does gaining physical access to a new rail station improve the labor-market outcomes of nearby residents? I focus on two outcomes---worker flows and income---and argue that a new station can affect both through three mechanisms: enabling workers to more reliably retain existing jobs, expanding their effective job-search radius across the metro area, and widening referral networks by connecting residents to a geographically broader set of social ties. Together, these suggest that transit access is not merely a convenience amenity but a potential driver of upward economic mobility---one distributed very unequally across urban space.
 
\section{Identification Strategy}
 
Identifying transit's causal effect is nontrivial. Station locations are not randomly chosen---municipalities tend to build in dense or well-resourced neighborhoods already experiencing growth. Residents also sort endogenously near planned stations, and long planning horizons create anticipatory effects in which firms and households adjust before a station ever opens.
 
I address these challenges by exploiting quasi-random variation in the \textit{timing} of station openings generated by construction delays. Conditional on a station having been selected for a neighborhood, I argue that delay length is as good as random---driven by engineering complications, permitting bottlenecks, and supply chain disruptions unrelated to local labor-market trends. By varying only \textit{when} a station opens rather than \textit{where}, I implicitly identify the average treatment effect on the treated, sidestepping station-placement endogeneity without modeling cities' location choices.
 
I embed this design into a stacked difference-in-differences framework, forming sub-experiments for each opening cohort that compare nearby tracts to tracts remaining untreated throughout the event window. Stacking these sub-experiments yields a balanced event-study design robust to heterogeneous treatment timing. Using planned opening dates from Environmental Impact Statements as my baseline time index also allows me to directly test for anticipatory adjustment and examine how excess delays distort residents' responses.
 
\section{Data}
 
A central contribution is a novel dataset of U.S.\ rail transit projects since 2000, linking planned opening dates, realized opening dates, and station geographies across projects in New York City, Washington D.C., San Francisco, Los Angeles, Seattle, Boston, Miami, and Honolulu. Planned dates are drawn from Draft Environmental Impact Statements archived in the EPA's EIS database---ensuring they reflect genuine ex-ante projections rather than post-delay revisions. Labor-market outcomes come from the Census Bureau's LODES Origin-Destination Employment Statistics, which records job flows between residential and workplace Census tracts. Transit exposure is measured using walking-time isochrones computed via the \texttt{r5r} package \parencite{r5r}, which routes tract-level travel times through OpenStreetMap pedestrian networks, replacing cruder Euclidean distance measures and accounting for real-world barriers to station access.
 
\section{Contribution}
 
This paper provides new causal evidence on how public transit shapes labor-market mobility, with direct implications for policy debates about transit investment and the geographic concentration of economic opportunity.
 \newpage
\printbibliography
 
\end{document}
